\begin{table*}[t]
\centering
\caption{Formal Hypothesis and Sub-Research Question Decisions based on 84-Run Pilot Benchmark.}
\label{tab:hypotheses}
\small
\begin{tabular}{lp{5.5cm}cp{7.5cm}}
\hline
\textbf{ID} & \textbf{Pre-Registered Claim / Question} & \textbf{Verdict} & \textbf{Primary Empirical Evidence} \\
\hline
H1 & H1: Asymmetry-Driven Decision Boundary Displacement & \textbf{PARTIALLY SUPPORTED} & Empirical evidence demonstrates that asymmetric noise ($\|T - T^\top\|_F > 0$) creates severe directional degradation in uncorrected empirical risk minimization (CE acc = 81.62\%), which is substantially mitigated by exact matrix inversion (Forward True $T$ acc = 87.80\%, $\Delta = +6.19\%$, $t = 22.77$, $p = 0.0019$, Holm $p = 0.0442$, Cohen's $d = 13.15$). Under symmetric noise, GCE matches or exceeds True $T$ without matrix inversion. Partially supported pending linear synthetic hyperplane angle $\angle(w^*, \tilde{w})$ measurement. \\
H2 & H2: Parametric Generalisation Window \& Capacity & \textbf{PARTIALLY SUPPORTED} & On PreActResNet-18 ($p/N \approx 240$), uncorrected CE exhibits memorization decay: validation accuracy drops by 2.99\% from peak in Symmetric 0.5 and by 2.29\% in Asymmetric 0.4. Forward Correction (True $T$) substantially arrests this drop to 1.00\% (Sym 0.5) and 0.27\% (Asym 0.4), while GCE suppresses it to 0.11\%. Partially supported pending variable capacity scaling grid ($p/N \in \{10^2, 10^4, 10^7\}$). \\
H3 & H3: Divergent Miscalibration Profiles \& Corrupted Validation Recovery & \textbf{SUPPORTED} & Empirically supported across all evaluated regimes. Post-hoc Temperature Scaling fitted on corrupted validation sets consistently underperforms clean-validation tuning across noisy regimes. Under Symmetric 0.2, CE test ECE under clean TS is 0.0779, but under corrupted TS it surges to 0.2442 ($\Delta = +0.1663$, $p = 0.0276$). Under Symmetric 0.5, True $T$ clean TS achieves ECE 0.0169 vs corrupted TS 0.1061 ($\Delta = +0.0892$, $p = 0.0008$). SCE exhibits severe underconfidence ($T^* > 2.0$, raw ECE up to 0.338). \\
H4 & H4: Controlled Confounder Audit on Human Noise & \textbf{OUT OF SCOPE} & Formally re-scoped as out of scope / future work. The benchmark is strictly focused on class-conditional synthetic noise. Real-world instance-dependent human noise (CIFAR-10N) is deferred to future work. \\
SRQ1 & SRQ 1: Theoretical Excess Risk Bounds vs Matrix Error \& Condition Number & \textbf{EMPIRICALLY QUANTIFIED} & Empirical results validate Proposition 2: condition number $\kappa(\hat{T})$ exhibits extreme sensitivity under Confident Learning in Asymmetric noise ($\kappa = 34.66 \pm 3.47$, $\|\hat{T}^{-1}\|_2 = 33.13$), driving test accuracy down from 87.80\% to 82.58\% and inflating ECE to 0.1193. \\
SRQ2 & SRQ 2: Calibration Recovery on Corrupted Validation Sets & \textbf{EMPIRICALLY RESOLVED} & Tuning Temperature Scaling on corrupted validation sets produces significant calibration penalties: up to +16.63 percentage points in ECE (CE in Sym 0.2) and +8.92 percentage points (True $T$ in Sym 0.5). In contrast, in-training robust losses such as GCE maintain robust uncalibrated test ECE (0.0750 in Sym 0.2, 0.0909 in Sym 0.5) without requiring validation sets. \\
SRQ3 & SRQ 3: Parametric Generalisation Window & \textbf{PARTIALLY RESOLVED} & Under Asymmetric 0.4 noise on PreActResNet-18, the generalization peak occurs early (epoch $22.3 \pm 2.5$), after which CE suffers a 2.29\% accuracy drop due to memorization. Forward correction arrests this degradation (0.27\% drop). Full resolution across varying $p/N$ ratios awaits the multi-capacity benchmark. \\
SRQ4 & SRQ 4: Controlled Human Noise Transfer Gap & \textbf{OUT OF SCOPE} & Re-scoped as future work; benchmark scope is strictly class-conditional synthetic noise on CIFAR-10. \\
\hline
\end{tabular}
\end{table*}
